\subsection{Continuation along the nodal locus}\label{ssec:componentwise}

Retain the hypotheses of Theorem~\ref{thm:main}, and write
\[
 P=X_{\mathcal T},\qquad L=-K_P,\qquad
 V=H^0(P,L)\setminus\{0\},\qquad Y_f=\{f=0\}\subset P.
\]
Thus $P$ is a smooth projective toric fourfold and $L$ is nef, big and
basepoint-free. Its sections are indexed by $\Delta\cap N$.
For any fixed $t>0$, the change of coefficients
$a_n=c_nt^{H(n)}$ identifies this family with the coefficient family in
Section~\ref{ssec:mirrorfamily}. We use all of $V$, including sections with some zero coefficients.
We do not take a quotient by automorphisms of $P$.

Put $\nu=|Q|$ and $r=\nu-\operatorname{rank}K$. Let $S_\nu\subset V$ consist of reduced
hypersurfaces with exactly $\nu$ distinct ordinary double points and no
other singularities. We give this locally closed locus its reduced
analytic structure; below we also identify its local defining ideal.

A marking labels the nodes by $Q$. An orientation marking chooses a
generator of $H_3(L_q;\Z)$ for each link $L_q$ of a node. These choices
form a finite covering of $S_\nu$ once its local manifold structure is
established. The chosen generator induces an orientation of the vanishing
sphere in any sufficiently small smooth deformation.
A relation has \emph{full support} if every coefficient is nonzero.

\begin{proposition}[componentwise continuation]\label{prop:componentwise}
Let $S$ be a connected component of $S_\nu$ containing a nodal member
constructed in Theorem~\ref{thm:main}(3). Choose at that member the marking
and orientations obtained from the near-limit construction, and let
$\widetilde S$ be the connected component of the marking and orientation
cover containing that choice. Then:
\begin{enumerate}
\item $S$ is smooth of codimension $r$ in $V$.
\item For $s\in\widetilde S$ and any sufficiently nearby smooth member
  $Y_b$, let $\delta_q(b)$ be its local vanishing spheres. The integral map
  \[
   d_b:\Z^Q\longrightarrow H_3(Y_b;\Z),\qquad
   e_q\longmapsto[\delta_q(b)]
  \]
  satisfies
  \[
   \ker d_b=d_b^{-1}\bigl(\Tors H_3(Y_b;\Z)\bigr)=K.
  \]
  In particular, its image is free of rank $r$ and meets the ambient
  torsion subgroup only in zero.
\item Every $k\in K$ bounds an integral singular $4$-chain in $Y_b$.
  Such a chain can be obtained by transport of a bounding chain supplied
  by the near-limit tropical construction, with the usual prism
  corrections for choices of fundamental cycles.
\item The spheres have pairwise disjoint Lagrangian representatives for
  the restriction of a fixed toric K\"ahler form. The equivalence of
  smoothability of $X'_\sigma$ with a full-support relation among these
  spheres, and the symplectic-transition conclusion of Corollary~\ref{cor:smooth}, hold at every such $Y_b$.
\end{enumerate}
\end{proposition}

On $S$ itself, the relation assertion is an equality of sub-local systems
of the signed permutation local system generated by the nodes. It does
not give a path-independent marking of individual nodes. Monodromy on
the whole of $H_3(Y_b)$ need not be a signed permutation. The signed
permutation assertion concerns the local generators only.

The proof uses the existing near-limit comparison as an input.
Lemma~\ref{lem:component-smoothness} proves that $S_\nu$ is smooth at every point,
including components not meeting the construction, and identifies its
local defining ideal through independent vanishing-cycle periods.
The comparison with $K$ then continues along each reached component.
Proposition~\ref{prop:x9different} shows that there can be components
which do not meet the construction. The argument does not construct a projective small resolution of
the mirror. Transported chains need not admit tropical descriptions at
their endpoints.

More precisely, the imported near-limit results provide a member
$f_0\in S_\nu$, a coefficient neighbourhood $U_0$ in which the only
possible singularities are the marked Morse points, and a connected
smooth complement $U_0^\circ$. At a known smooth member of this
complement, the relation lattice and the preimage of ambient torsion
are both $K$, and every element of $K$ has a tropical bounding chain.
The original resolution-side criterion identifies smoothability of
$X'_\sigma$ with the existence of a full-support vector in $K$.
These are the inputs from Theorem~\ref{thm:main}, Lemma~\ref{lem:U},
Proposition~\ref{prop:relationsU} and Corollary~\ref{cor:smooth}.

\subsubsection*{Completeness of polynomial deformations}

For a face $F$ let $\ell^*(F)$ denote its number of relative interior
lattice points. The correction term for the Newton polytope $\Delta$ is
\begin{equation}\label{eq:component-correction}
 c(\Delta)=\sum_{\dim G=2}\ell^*(G)\ell^*(G^\circ)=0.
\end{equation}
Indeed, the triangles and squares have no interior lattice points, while
the dual edges of the pentagons and hexagons have lattice length one.
The dual edge length of a triangular face is immaterial here.

\begin{lemma}\label{lem:component-complete}
For every reduced anticanonical hypersurface $Y\subset P$,
$H^1(Y,T_P|_Y)=0$. If $Y$ has only nodes, the Kodaira--Spencer map
\[
 H^0(P,L)\longrightarrow\operatorname{Ext}^1_Y(\Omega_Y,\mathcal O_Y)
\]
of the coefficient family is surjective.
\end{lemma}

\begin{proof}
The cohomology of a smooth projective toric variety is generated by
algebraic cycle classes, so $H^j(P,\Omega_P^i)=0$ for $i\ne j$.
Also $H^j(P,\mathcal O_P)=0$ for $j>0$, and nef line bundles have no
higher cohomology. These are the standard toric cohomology results;
see \cite[Thms.~9.2.3, 9.4.11]{CLS}. Since
$T_P\otimes K_P\simeq\Omega_P^3$, the restriction sequence gives
\begin{equation}\label{eq:component-ambientrestriction}
 H^i(P,T_P)\simeq H^i(Y,T_P|_Y),\qquad i=0,1.
\end{equation}
This holds for every anticanonical Cartier divisor, independently of its
defining coefficients.

We compute $h^1(P,T_P)$ using a smooth $\Delta$-regular member $Z$.
For any anticanonical member $Y$, adjunction gives
$\omega_Y\simeq\mathcal O_Y$. The sequence
$0\to K_P\to\mathcal O_P\to\mathcal O_Y\to0$, Serre duality on
the fourfold $P$, and the preceding toric vanishing give
\[
 H^0(Y,\mathcal O_Y)=\C,\qquad
 H^1(Y,\mathcal O_Y)=H^2(Y,\mathcal O_Y)=0.
\]
In particular, $Y$ is connected. If $Y$ has only nodes, it is normal
and hence irreducible; thus the nodal members are projective
Calabi--Yau threefolds in the usual sense. For the smooth member $Z$
we obtain $K_Z\simeq\mathcal O_Z$ and $H^2(\mathcal O_Z)=0$.
Hence $H^0(T_Z)=H^0(\Omega_Z^2)=0$ and $h^1(T_Z)=h^{2,1}(Z)$.
The restriction sequence of $L$ gives
\[
 h^0(Z,L|_Z)=\ell(\Delta)-1,\qquad H^1(Z,L|_Z)=0.
\]
Taking cohomology in the normal sequence of $Z\subset P$ and using
\eqref{eq:component-ambientrestriction} yields
\begin{equation}\label{eq:component-h1ambient}
 h^1(P,T_P)=h^{2,1}(Z)-\ell(\Delta)+1+h^0(P,T_P).
\end{equation}
The toric Euler sequence \cite[Thm.~8.1.6]{CLS}, or equivalently the root
description of toric vector fields, gives
\begin{equation}\label{eq:component-roots}
 h^0(P,T_P)=4+\sum_{\dim F=3}\ell^*(F).
\end{equation}
For completeness, a nonconstant monomial section of a ray divisor
$D_\rho$ is indexed by a lattice point $m$ with
$\langle m,v_\rho\rangle=-1$ and
$\langle m,v_{\rho'}\rangle\ge0$ for every other ray.
Such a $v_\rho$ must be a vertex of $\Delta^\circ$, and $m$ is an
interior lattice point of its dual facet in $\Delta$. Conversely every
such point satisfies these inequalities for all the boundary rays:
its pairing can equal $-1$ only at that vertex, and the other integral
pairings are nonnegative. Summing $h^0(\mathcal O_P(D_\rho))$ and
subtracting the rank of $\operatorname{Pic}(P)$ in the Euler sequence
proves \eqref{eq:component-roots}.

Batyrev's formula \cite[Thm.~4.3.7]{Batyrev} is
\[
 h^{2,1}(Z)=\ell(\Delta)-5-
 \sum_{\dim F=3}\ell^*(F)+c(\Delta).
\]
Equations \eqref{eq:component-correction}--\eqref{eq:component-roots} therefore give
$H^1(P,T_P)=0$, proving the first assertion.

For a reduced nodal hypersurface the conormal sequence
\[
 0\longrightarrow L^{-1}|_Y\xrightarrow{df}
 \Omega_P|_Y\longrightarrow\Omega_Y\longrightarrow0
\]
is exact on the left as well: $df$ is nonzero at the generic smooth
point of every component, and $\mathcal O_Y$ is reduced.
Its global Ext sequence contains
\[
 H^0(Y,L|_Y)\longrightarrow\operatorname{Ext}^1_Y(\Omega_Y,\mathcal O_Y)
 \longrightarrow H^1(Y,T_P|_Y)=0.
\]
Finally $H^0(P,L)\to H^0(Y,L|_Y)$ is surjective because
$H^1(P,\mathcal O_P)=0$. The composite is the claimed
Kodaira--Spencer map.
\end{proof}

\subsubsection*{The rank of the node-smoothing map}

Fix $f\in S_\nu$ and write $Y=Y_f$. Choose mutually disjoint small
conical neighbourhoods of its nodes, and let $W$ be their compact
exterior. Its oriented boundary is a disjoint union
$L_\partial=\coprod_{q\in Q}L_q$ with $L_q\simeq S^2\times S^3$.
Write the integral inclusion maps as
\begin{equation}\label{eq:component-ab}
 a:H_3(L_\partial;\Z)\longrightarrow H_3(W;\Z),\qquad
 b:H_2(L_\partial;\Z)\longrightarrow H_2(W;\Z).
\end{equation}
The orientation marking identifies the source of $a$ with $\Z^Q$.
Use the boundary intersection pairing to choose the dual basis in the
source of $b$. In this section ranks are computed over $\Q$.

\begin{lemma}\label{lem:component-topology}
One has $\operatorname{rank} a+\operatorname{rank} b=\nu$. If $Y_u$ is a sufficiently small
smooth deformation of $Y$, then
\begin{equation}\label{eq:component-integralmv}
 0\longrightarrow H_3(W;\Z)\xrightarrow{j}H_3(Y_u;\Z)
 \longrightarrow\ker b\longrightarrow0.
\end{equation}
Under suitable compatible signs, its vanishing-cycle map is $d_u=j a$.
In particular,
\begin{equation}\label{eq:component-torsioncomparison}
 \ker d_u=\ker a,\qquad
 d_u^{-1}(\Tors H_3(Y_u;\Z))=a^{-1}(\Tors H_3(W;\Z)).
\end{equation}
\end{lemma}

\begin{proof}
The boundary map $H_3(W,L_\partial;\Q)\to H_2(L_\partial;\Q)$
is adjoint to $a$ by Poincar\'e--Lefschetz duality and the intersection
pairing on $L_\partial$. Its image is $\ker b$ by the exact sequence of
the pair. Thus $\dim\ker b=\operatorname{rank} a$, giving the rank identity.

Replace each conical neighbourhood by its Milnor fibre $F_q$.
The fibre is a disc cotangent bundle of $S^3$, and
$H_3(L_q;\Z)\to H_3(F_q;\Z)$ is an isomorphism. One way to see the
integral assertion is to trivialise $T^*S^3$: the map is the projection
$S^3\times S^2\to S^3$. Also $H_2(F_q;\Z)=0$.
Mayer--Vietoris for $Y_u=W\cup\coprod_qF_q$ therefore has the segment
\[
 \Z^Q\xrightarrow{(a,-1)}H_3(W;\Z)\oplus\Z^Q
 \longrightarrow H_3(Y_u;\Z)
 \longrightarrow\Z^Q\xrightarrow{b}H_2(W;\Z).
\]
The cokernel of $(a,-1)$ identifies with $H_3(W;\Z)$ by
$(w,v)\mapsto w+a(v)$. This proves \eqref{eq:component-integralmv} and
$d_u=ja$ after orienting the local core spheres compatibly.
The group $\ker b$ is free, being a subgroup of $\Z^Q$.
Consequently $j$ identifies the torsion subgroup of its source with that
of its target, giving \eqref{eq:component-torsioncomparison}.
\end{proof}

Let
\[
 \lambda:\operatorname{Ext}^1_Y(\Omega_Y,\mathcal O_Y)
       \longrightarrow\bigoplus_{q\in Q}T^1_q\simeq\C^Q
\]
be the global-to-local node-smoothing map. The choice of coordinates on
the one-dimensional $T^1_q$ affects it by nonzero diagonal factors only.

\begin{lemma}\label{lem:component-rank}
The rank of $\lambda$ equals the rank of the span of the vanishing
cycles in a nearby smoothing. Both equal $\operatorname{rank} a$.
\end{lemma}

\begin{proof}
Let $\pi:\overline Y\to Y$ be the blow-up of all nodes. This is a
smooth projective threefold, with exceptional quadrics
$E_q\simeq\mathbb P^1\times\mathbb P^1$ of normal bundle
$\mathcal O_{E_q}(-1,-1)$. Denote their ruling classes in
$H_2(\overline Y;\Q)$ by $A_q,B_q$, and put $D_q=A_q-B_q$.
The node-smoothing sequence of Rollenske--Thomas
\cite[\S2 and (3.2)]{RT09} identifies the image of $\lambda$ with the
kernel of the map sending the $q$th generator to $D_q$ in
\[
 H^{2,2}(\overline Y)\big/\langle A_q+B_q:q\in Q\rangle.
\]
Here curve classes are identified with degree-four cohomology classes.
All hypotheses hold: $Y$ has trivial dualizing sheaf by adjunction,
has only ordinary double points, and its blow-up is K\"ahler because it
is projective. No small resolution is used.
The span of the $D_q$ has zero intersection with the span of the
$A_q+B_q$: pairing with each exceptional divisor gives zero on every
$D_q$ and $-2$ on its own sum, with zero on the other sums.
Hodge decomposition on $\overline Y$ therefore shows
\begin{equation}\label{eq:component-rt-rank}
 \operatorname{rank}\lambda=\nu-\dim_{\Q}\langle D_q:q\in Q\rangle.
\end{equation}

We identify the last dimension without assuming that a small resolution
is K\"ahler. Topologically $\overline Y$ is obtained from $W$ by
attaching the disc bundles of $\mathcal O_{E_q}(-1,-1)$. The map from
$H_2(L_q;\Q)$ to $H_2(E_q;\Q)$ is injective with image
$\Q(A_q-B_q)$, by the circle-bundle sequence. After compatible choices
of sign, Mayer--Vietoris gives
\[
 H_2(\overline Y;\Q)=
 \frac{H_2(W;\Q)\oplus\bigoplus_qH_2(E_q;\Q)}
 {\{(b(v),-D(v)):v\in\Q^Q\}}.
\]
The map $D:\Q^Q\to\bigoplus_qH_2(E_q;\Q)$ is injective.
It follows that $\sum v_qD_q$ vanishes in $H_2(\overline Y;\Q)$
exactly when $b(v)=0$. Thus the dimension in \eqref{eq:component-rt-rank}
is $\operatorname{rank} b$. Lemma~\ref{lem:component-topology} now gives
$\operatorname{rank}\lambda=\nu-\operatorname{rank} b=\operatorname{rank} a=\operatorname{rank} d_u$.
\end{proof}

\subsubsection*{Local equations of the nodal locus}

Choose a sufficiently small coefficient neighbourhood $U\subset V$ of
$f\in S_\nu$. The holomorphic Morse lemma with parameters gives one
holomorphic critical point near each node and a critical-value function
$\varepsilon_q:U\to\C$. In suitable local coordinates the hypersurface
equation is
\begin{equation}\label{eq:component-morse}
 z_1^2+z_2^2+z_3^2+z_4^2=\varepsilon_q(u).
\end{equation}
Multiplying the original local equation by a unit changes the
coordinates and the critical value, but not its vanishing locus or the
rank of their differentials at $f$. Compactness of the complement of
the node neighbourhoods ensures, after shrinking $U$, that these are
all possible singular points. Each is a node whenever its critical
value vanishes.

The derivative $d\varepsilon_q(f)$ is, up to a nonzero factor,
evaluation of a section at the node. Basepoint-freeness of $L$ makes
each such derivative nonzero. A tangent direction outside their
finitely many kernels gives a simultaneous smoothing. The set
\[
 U^\circ=U\setminus\bigcup_{q\in Q}\{\varepsilon_q=0\}
\]
is nonempty and connected after choosing a small ball for $U$.
Over it, the oriented local vanishing cycles are single-valued as
homology classes. A loop of the local parameter in \eqref{eq:component-morse}
acts by the antipodal map on $S^3$, of degree $+1$; the disjoint local
charts give the same conclusion for arbitrary loops in $U^\circ$.

Choose the relative holomorphic volume form $\Omega_u$ by adjunction
from the defining section and a fixed trivialization of $K_P\otimes L$.
The periods
\[
 p_q(u)=\int_{\delta_q(u)}\Omega_u
\]
extend holomorphically over $U$, and
\begin{equation}\label{eq:component-periodunit}
 p_q=h_q\varepsilon_q,\qquad h_q(f)\ne0.
\end{equation}
Indeed, in the Morse chart the residue is the residue of
$g(z,u)\,dz_1\cdots dz_4/(\sum z_i^2-\varepsilon_q(u))$ with
$g(0,f)\ne0$. Scaling $z=\sqrt{\varepsilon_q}\,x$ contributes a factor
$\varepsilon_q$. The integrals of the odd-degree Taylor terms of $g$
vanish under the antipodal map, and the even-degree terms give a
convergent power series in $\varepsilon_q$. The leading integral is a
nonzero constant times $g(0,f)$. Shrink $U$ so that every $h_q$ is a unit.

\begin{lemma}\label{lem:component-smoothness}
The locus $S_\nu$ is smooth near $f$, with codimension equal to
$c=\operatorname{rank} a$. Its local ideal is $(\varepsilon_q:q\in Q)$ and is
generated by $c$ independent periods.
\end{lemma}

\begin{proof}
Lemma~\ref{lem:component-complete} and Lemma~\ref{lem:component-rank} imply
$\operatorname{rank} d\varepsilon(f)=c$: the differential of the critical-value map
is the composite of the Kodaira--Spencer map with $\lambda$, followed
by nonzero diagonal factors. Equation~\eqref{eq:component-periodunit} gives
$\operatorname{rank} dp(f)=c$ as well.

Fix $u\in U^\circ$ and put
$R_{\Q}=\ker(d_u\otimes\Q)\subset\Q^Q$. This space has dimension
$\nu-c$ and is independent of $u$ by parallel transport of the local
cycles. For $k\in R_{\Q}$ one has
\begin{equation}\label{eq:component-periodrelations}
 \sum_q k_qp_q=0
\end{equation}
on $U^\circ$, hence on $U$ by holomorphic extension. Its complexification
is therefore contained in the space of row relations of $dp(f)$.
Both spaces have dimension $\nu-c$, so they are equal.
Choose an index set $I\subset Q$ of size $c$ whose rows in $dp(f)$ are
independent. For each $q\notin I$, the corresponding row dependence
belongs to $R_{\Q}\otimes\C$; equation~\eqref{eq:component-periodrelations}
therefore expresses $p_q$ as a constant linear combination of the
$p_i$, $i\in I$, throughout $U$.
Thus
\[
 (\varepsilon_q:q\in Q)=(p_q:q\in Q)=(p_i:i\in I).
\]
The last generators have independent differentials at $f$.
The implicit function theorem proves the assertion, including
reducedness of this local zero scheme. By the choice of $U$, its
points are exactly those with all $\nu$ nodes present and no other
singularities, so it is $S_\nu\cap U$.
\end{proof}

The smoothness statement follows from completeness of polynomial
deformations, the rank calculation, and the period identities.
It applies at every point of $S_\nu$, including components for which
no comparison with $K$ has yet been obtained.

\subsubsection*{Integral continuation and the symplectic consequence}

\begin{proof}[Proof of Proposition~\ref{prop:componentwise}]
Lemma~\ref{lem:component-smoothness} gives $S_\nu$ its local manifold structure.
The universal set of nodes is a finite unramified covering of degree
$\nu$. Locally along this locus, the parametrized Morse coordinates
trivialize the singularities and their links. The exterior of small
node neighbourhoods is a proper smooth family of compact manifolds
with boundary. Ehresmann's theorem, also for the boundary, supplies
local differentiable trivializations. Choices of sufficiently small
neighbourhoods agree up to isotopy and collar identifications. Equivalently,
the conical collars identify $H_3(W;\Z)$ with
$H_3(Y_{\mathrm{reg}};\Z)$, where $Y_{\mathrm{reg}}$ is the smooth locus
of the nodal hypersurface.
Consequently $H_3(L_\partial;\Z)$ and $H_3(W;\Z)$ are local systems,
with the inclusion-induced morphism $a$ in \eqref{eq:component-ab}.

Ordering the nodes and orienting their rank-one link groups kills the
finite signed permutation monodromy of the source. Its connected
covering $\widetilde S$ through the initial marking surjects onto $S$.
On this cover the source is the constant group $\Z^Q$.
Both $\ker a$ and $a^{-1}(\Tors H_3(W;\Z))$ are constant subgroups
of that source: transport in the target preserves the inclusion map
and its torsion subgroup.
At the near-limit starting point, Proposition~\ref{prop:relationsU} and Theorem~\ref{thm:main}, together with \eqref{eq:component-torsioncomparison}, identify both
subgroups with $K$. They therefore equal $K$ throughout
$\widetilde S$. Equation~\eqref{eq:component-torsioncomparison} proves part~(2)
for every sufficiently nearby smooth deformation of each member.
The rank $c$ in Lemma~\ref{lem:component-smoothness} is now $r$ throughout $S$,
proving part~(1).

Here is also an integral chain description. Represent the marked link
generators by cycles $\beta_q$. Since $a(k)=0$ for $k\in K$, there is
an integral $4$-chain $C_W$ in $W$ with
$\partial C_W=\sum_q k_q\beta_q$.
In the local Milnor fibre, $\beta_q$ and the vanishing sphere represent
the same generator of $H_3(F_q;\Z)$, so an integral $4$-chain $C_q$
has boundary $[\delta_q]-\beta_q$. Then
\[
 \partial\left(C_W+\sum_qk_qC_q\right)
       =\sum_qk_q[\delta_q].
\]
To retain a specific chain from the tropical construction, take a path
in $\widetilde S$ from the initial member to the desired one.
First transport that chain inside $U_0^\circ$ from its constructed
smooth member to a sufficiently small smoothing of $f_0$;
connectedness of $U_0^\circ$ and the flatness of the marked spheres
allow this. This initial step uses the original neighbourhood $U_0$,
which need not equal the smaller neighbourhood of $f_0$ used below.
Cover this compact path by finitely many coefficient neighbourhoods
as above and subdivide it accordingly. At consecutive overlaps choose
smooth members arbitrarily near the intervening nodal member.
Within each neighbourhood the complement of the discriminant is
path connected, so these choices can be joined to give a path of
smooth hypersurfaces from the initial smoothing to the target one.
Along each piece the local cycles identify with the marked link
generators. Differentiable parallel transport carries the initial
integral chain to the final fibre. Local sphere isotopies and prism
chains adjust its boundary to the specified fundamental cycles.
This proves part~(3); no uniqueness of the resulting chain is asserted.

Fix a toric K\"ahler form on $P$. At a nodal member and a sufficiently
small simultaneous smoothing, the local holomorphic Morse functions
have Lefschetz thimbles contained in mutually disjoint node
neighbourhoods. Their intersections with the smooth fibre are
embedded Lagrangian $3$-spheres, smoothly isotopic to the local Milnor
spheres. The containment follows by taking the critical values small
compared with a fixed ball in each Morse chart; on that ball the
K\"ahler metric is uniformly comparable with the Euclidean metric.
This is precisely the local argument of Lemma~\ref{lem:L} and does
not use its near-limit estimates. A path in $U^\circ$ transports the
spheres to any other smooth member of the chosen neighbourhood.
The restriction of the fixed ambient K\"ahler form gives symplectic
parallel transport on this proper smooth family, preserving
disjointness and the vanishing homology classes.

The original resolution-side criterion in Corollary~\ref{cor:smooth} says
that $X'_\sigma$ is smoothable exactly when $K$ contains an integral
vector with every coordinate nonzero. Part~(2) transfers this criterion
to the mirror spheres. Smith--Thomas--Yau \cite[Thm.~2.9]{STY} then
provides a symplectic conifold surgery in one of the choices of local
small resolutions, with symplectic exceptional curves. This proves
part~(4).
\end{proof}

\begin{remark}[Scope and boundary cases]
The statement concerns sufficiently small smoothings of members of a
reached nodal component. It does not assert a canonical collection of
spheres on every smooth member of the whole coefficient space.
One may transport a configuration farther along a chosen path of
smooth members, but that is additional path-dependent data.
Scalar multiplication acts freely on $V$ and preserves $S_\nu$;
the same smoothness and codimension statement holds in $\mathbb P H^0(P,L)$.
A coarse quotient by ambient automorphisms is a different question.
If $\nu=0$, the locus is open and the relation assertions are empty.
Colliding nodes, additional singularities and nonreduced hypersurfaces
lie outside the stated locus.
\end{remark}
